\begin{figure}[ht]
  \centering
  % NIENTE \resizebox: ottimizzato orizzontalmente per far rientrare la grafica 
  % nei margini e lasciare che i font \footnotesize si renderizzino alla loro grandezza reale.
  \begin{tikzpicture}[
      % Distanza orizzontale compatta per accorciare le frecce laterali
      node distance=3.5ex and 4.5em,
      %
      % Stili di base
      base/.style={draw, rounded corners=1.5pt, inner sep=5pt, align=center, font=\small, minimum height=4.5ex},
      %
      % Stili Fasi (Colonna Sinistra)
      phase-startup/.style={base, fill=blue!7, draw=blue!40, minimum width=13.5em},
      phase-exec/.style={base, fill=green!7, draw=green!45, minimum width=13.5em},
      %
      % Stili Errori (Colonna Destra)
      err-stop/.style={base, fill=red!6, draw=red!50, minimum width=12.5em},
      err-warn/.style={base, fill=orange!8, draw=orange!50, minimum width=12.5em},
      %
      % Terminali (Inizio/Fine)
      term/.style={draw=gray!50, fill=gray!15, ellipse, minimum height=3.5ex, inner sep=5pt, align=center, font=\footnotesize\ttfamily},
      %
      % Testi secondari sulle frecce (inner sep=3pt aggiunto per distaccarle dai fili)
      lbl/.style={font=\footnotesize\itshape, text=gray!70!black, inner sep=3pt},
      %
      % Frecce ortogonali
      arr-ok/.style={-latex, semithick, gray!70!black},
      arr-ko/.style={-latex, semithick, red!60!black},
      arr-warn/.style={-latex, semithick, dashed, orange!80!black},
      %
      % Sfondi per raggruppamento dinamico (Libreria fit)
      bg-startup/.style={draw=blue!30, fill=blue!3, dashed, rounded corners=3pt, inner sep=10pt},
      bg-exec/.style={draw=green!35, fill=green!2, dashed, rounded corners=3pt, inner sep=10pt}
    ]

    %% ── Terminale di avvio ─────────────────────────────────────────────────
    \node[term] (start) {apiguard run};

    %% ── Fasi 1-4 (zona startup, bloccanti) ─────────────────────────────────
    % Spazio a 8ex per creare il corridoio d'ingresso
    \node[phase-startup, below=8ex of start] (p1) {\textbf{Fase~1:} Inizializzazione\\{\footnotesize\texttt{config/loader.py}}};
    \node[phase-startup, below=of p1]    (p2) {\textbf{Fase~2:} OpenAPI Discovery\\{\footnotesize\texttt{discovery/openapi.py}}};
    \node[phase-startup, below=of p2]    (p3) {\textbf{Fase~3:} Costruzione Contesti\\{\footnotesize\texttt{engine.py}}};
    \node[phase-startup, below=of p3]    (p4) {\textbf{Fase~4:} Test Discovery e \gls{DAG}\\{\footnotesize\texttt{registry.py, dag.py}}};

    %% ── Fasi 5-7 (zona esecuzione, non bloccanti) ──────────────────────────
    % GAP ridotto: spazio a 9ex (lascia 5ex netti tra i due box di background)
    \node[phase-exec, below=9ex of p4] (p5) {\textbf{Fase~5:} Esecuzione (per test)\\{\footnotesize\texttt{engine.py}}};
    \node[phase-exec, below=of p5]   (p6) {\textbf{Fase~6:} Teardown \gls{LIFO}\\{\footnotesize\texttt{core/context.py}}};
    \node[phase-exec, below=of p6]   (p7) {\textbf{Fase~7:} Report\\{\footnotesize\texttt{report/renderer.py}}};

    %% ── Terminale di uscita ────────────────────────────────────────────────
    \node[term, below=5ex of p7] (end) {exit code: 0 / 1 / 2 / 10};

    %% ── Rami KO → STOP (fasi bloccanti) posizionati prima degli sfondi ─────
    \node[err-stop, right=4.5em of p1] (s1) {\texttt{ConfigurationError}\\{\footnotesize avvio interrotto}};
    \node[err-stop, right=4.5em of p2] (s2) {\texttt{OpenAPILoadError}\\{\footnotesize avvio interrotto}};
    \node[err-stop, right=4.5em of p4] (s4) {\texttt{DAGCycleError}\\{\footnotesize avvio interrotto}};

    %% ── Rami errore recuperabile (fasi 5, 6) ───────────────────────────────
    \node[err-warn, right=4.5em of p5] (w5) {\texttt{TestResult(ERROR)}\\{\footnotesize isolato, pipeline prosegue}};
    \node[err-warn, right=4.5em of p6] (w6) {\texttt{TeardownError}\\{\footnotesize \texttt{WARNING} strutturato}};

    %% ── Sfondi per le due zone ─────────────────────────────────────────────
    \begin{scope}[on background layer]
      \node[bg-startup, fit=(p1)(p2)(p3)(p4)] (grp-startup) {};
      \node[bg-exec, fit=(p5)(p6)(p7)] (grp-exec) {};
    \end{scope}

    %% ── ETICHETTE DI GRUPPO (Avvicinate all'asse e centrate in Y) ──────────
    % xshift=-0.6em le avvicina notevolmente alla freccia centrale
    % yshift=3ex solleva la prima etichetta al centro matematico del gap superiore.
    \node[font=\footnotesize\bfseries, text=blue!70!black, anchor=east]
    at ([xshift=-0.6em, yshift=3ex]grp-startup.north) {Fasi 1-4 (bloccanti)};

    % yshift=2.5ex solleva la seconda etichetta al centro matematico del gap centrale di 5ex netti.
    \node[font=\footnotesize\bfseries, text=green!70!black, anchor=east]
    at ([xshift=-0.6em, yshift=2.5ex]grp-exec.north) {Fasi 5-7 (non bloccanti)};

    %% ── Frecce asse principale ─────────────────────────────────────────────
    \draw[arr-ok] (start) -- (p1);
    \draw[arr-ok] (p1) -- node[midway, right, lbl] {OK} (p2);
    \draw[arr-ok] (p2) -- node[midway, right, lbl] {OK} (p3);
    \draw[arr-ok] (p3) -- node[midway, right, lbl] {OK} (p4);
    \draw[arr-ok] (p4) -- node[midway, right, lbl] {OK} (p5);
    \draw[arr-ok] (p5) -- (p6);
    \draw[arr-ok] (p6) -- node[midway, right, lbl] {OK} (p7);
    \draw[arr-ok] (p7) -- (end);

    %% ── Frecce Laterali (Centrate nello spazio vuoto ricalibrato a pos=0.55) 
    % Il valore 0.55 arretra i testi rispetto al 0.65 precedente, staccandoli dai riquadri di destra
    \draw[arr-ko] (p1.east) -- node[pos=0.55, above, lbl, text=red!70!black] {KO} (s1.west);
    \draw[arr-ko] (p2.east) -- node[pos=0.55, above, lbl, text=red!70!black] {KO} (s2.west);
    \draw[arr-ko] (p4.east) -- node[pos=0.55, above, lbl, text=red!70!black] {KO} (s4.west);

    \draw[arr-warn] (p5.east) -- node[pos=0.55, above, align=center, lbl, text=orange!80!black] {eccezione\\interna} (w5.west);
    \draw[arr-warn] (p6.east) -- node[pos=0.55, above, align=center, lbl, text=orange!80!black] {cleanup\\parziale} (w6.west);

  \end{tikzpicture}
  \caption{Flusso della pipeline a sette fasi (D4.P3, D4.P4): la zona di startup (Fasi~1-4) è bloccante, un errore in una qualsiasi di esse interrompe l'avvio; la zona di esecuzione (Fasi~5-7) è non bloccante, ogni malfunzionamento è recuperato localmente senza invalidare i risultati già prodotti. Frecce continue: percorso nominale; frecce tratteggiate: errori recuperati con \texttt{WARNING}.}
  \label{fig:pipeline-flowchart}
\end{figure}